\chapter{تصویر کلی علوم زبان و آواشناسی تولیدی}
\emergencystretch=4em
\newcommand{\chtwoFz}{F_{\text{\lr{0}}}}
\newcommand{\chtwoTz}{T_{\text{\lr{0}}}}

\section{مرور جلسهٔ قبل و تغییر رویکرد}

در جلسهٔ نخست مسئله را با رویکرد \term{بالا‌به‌پایین}{Top-Down} دیدیم: ابتدا تسک‌های زبان (مانند \eng{POS Tagging}، \eng{NER}، \eng{Language Modeling}) و گفتار (مانند \eng{STT}، \eng{TTS}، \eng{Speaker Diarization}) و سپس ترکیب آن‌ها در یک \eng{Pipeline}. از این جلسه مسیر را برعکس می‌کنیم و رویکرد \term{پایین‌به‌بالا}{Bottom-Up} را در پیش می‌گیریم:
\begin{enumerate}
  \item ابتدا خود \textbf{داده} را بشناسیم: از چه جنسی است و چگونه به وجود می‌آید؟
  \item سپس \term{بازنمایی}{Representation}: داده را چگونه به شکلی درآوریم که رایانه پردازش کند؟ (دست‌ساز، یادگرفته‌شده یا یادگرفته‌شده با خودنظارتی)
  \item در پایان، مدل‌های پیچیده‌تر و عمیق‌تر.
\end{enumerate}

جنس دادهٔ گفتاری و متنی یکسان نیست: متن دنباله‌ای از نمادهاست، اما گفتار یک پدیدهٔ فیزیکی است که در هوا منتشر می‌شود.

\begin{table}[htbp]
\centering
\caption{مقایسهٔ ماهیت داده در متن و گفتار}
\begin{tabular}{p{0.38\textwidth} p{0.38\textwidth}}
\toprule
\textbf{گفتار} (\eng{Speech}) & \textbf{متن} (\eng{Text}) \\
\midrule
ماهیت فیزیکی دارد & ماهیت نمادین دارد \\
به‌صورت موج صوتی منتقل می‌شود & از نمادها و نشانه‌ها تشکیل شده است \\
در زمان پیوسته است & در سطح توالی نمادها گسسته است \\
\bottomrule
\end{tabular}
\end{table}

بنابراین در این جلسه ابتدا روی ماهیت دادهٔ گفتاری تمرکز می‌کنیم: موج صوتی چیست، چگونه ایجاد می‌شود و انسان چگونه صحبت می‌کند.

\section{تصویر کلی: جایگاه پردازش گفتار در علوم زبان}

\term{زبان‌شناسی}{Linguistics} در کنار فیزیک، ریاضی، روان‌شناسی و جامعه‌شناسی از علوم مرتبط با موضوع ماست و شاخه‌های متعددی دارد:
\begin{description}
  \item[\term{آواشناسی}{Phonetics}] مطالعهٔ صداهای گفتاری انسان: تولید، ویژگی‌های فیزیکی و شنیده‌شدن. محور اصلی این جلسه و جلسات بعد.
  \item[\term{نحو}{Syntax}] قواعد چینش کلمات برای ساختن جمله؛ در \eng{NLP} بسیار با آن سروکار داریم.
  \item[\term{ریخت‌شناسی (صرف)}{Morphology}] ساختار درونی کلمات و تغییر شکل آن‌ها: صرف فعل «رفتن» (رفتم، رفتی، \ldots) یا واژه‌سازی «دانش» + «ـمند» $\leftarrow$ «دانشمند».
\end{description}
\noindent دو علم میان‌رشته‌ای مرتبط:
\begin{description}
  \item[\term{زبان‌شناسی رایانشی}{Computational Linguistics}] در تقاطع زبان‌شناسی، علوم کامپیوتر و هوش مصنوعی؛ مدل‌سازی محاسباتی زبان. درس ما، \term{پردازش زبان گفتاری}{Spoken Language Processing}، در همین چارچوب است.
  \item[\term{عصب‌زبان‌شناسی}{Neurolinguistics}] در تقاطع زبان‌شناسی و علوم اعصاب؛ سازوکارهای عصبی مغز برای پردازش زبان.
\end{description}

\begin{exambox}{نگاه مهندسی به زبان‌شناسی}
هدف ما به‌عنوان مهندس هوش مصنوعی ارائهٔ نظریهٔ تازه در زبان‌شناسی یا کشف دوبارهٔ دسته‌بندی‌های آن نیست؛ باید مفاهیم و دسته‌بندی‌های موردنیاز را بشناسیم و در ساخت مدل‌های محاسباتی از آن‌ها بهره بگیریم.
\end{exambox}

\section{نظریهٔ زایشی چامسکی (۱۹۵۷)}

\begin{itemize}
  \item \textbf{دیدگاه تقلیدی (قدیمی):} کودک جمله‌ها را می‌شنود، حفظ می‌کند و تکرار می‌کند.
  \item \textbf{نقد چامسکی:} کودکان مرتب \textbf{جمله‌های جدیدی} می‌سازند که هرگز نشنیده‌اند. پس توانایی زبانی بیش از حفظ است و در ذهن سازوکاری برای «تولید» (\eng{generate}) جمله‌ها وجود دارد؛ این ایده به شکل‌گیری علوم شناختی کمک کرد.
  \item \textbf{زبان خطی نیست، \term{سلسله‌مراتبی}{Hierarchical} و درختی است:} کلمات در گروه‌ها (اسمی، فعلی) سازمان می‌یابند و گروه‌ها در هم تو در تو می‌شوند (شکل \ref{fig:ch02-tree}).
\end{itemize}

\begin{figure}[!ht]
\centering
\begin{tikzpicture}[rtl,
  level distance=11mm,
  every node/.style={font=\small},
  cat/.style={font=\small\itshape}]
\node (S) at (0,0) {\lr{S}};
\node (NP1) at (1.8,-1.2) {\lr{NP}};
\node (VP) at (-1.8,-1.2) {\lr{VP}};
\node (N1) at (1.8,-2.4) {\lr{N}};
\node (w1) at (1.8,-3.4) {کودک};
\node (NP2) at (0,-2.4) {\lr{NP}};
\node (V) at (-3,-2.4) {\lr{V}};
\node (w3) at (-3,-3.4) {خورد};
\node (N2) at (0.8,-3.4) {\lr{N}};
\node (P) at (-0.8,-3.4) {\lr{P}};
\node (w2) at (0.8,-4.4) {سیب};
\node (w2b) at (-0.8,-4.4) {را};
\draw (S)--(NP1); \draw (S)--(VP);
\draw (NP1)--(N1); \draw (N1)--(w1);
\draw (VP)--(NP2); \draw (VP)--(V); \draw (V)--(w3);
\draw (NP2)--(N2); \draw (NP2)--(P); \draw (N2)--(w2); \draw (P)--(w2b);
\end{tikzpicture}
\caption{درخت ساختار سلسله‌مراتبی برای «کودک سیب را خورد». \lr{S}: جمله، \lr{NP}: گروه اسمی، \lr{VP}: گروه فعلی، \lr{N}: اسم، \lr{V}: فعل، \lr{P}: نشانهٔ «را». (شکل افزوده)}
\label{fig:ch02-tree}
\end{figure}

\paragraph{دستور جهانی در برابر یادگیری آماری.} پرسش اصلی: چه مقدار از توانایی زبانی از ابتدا در ذهن هست و چه مقدار از تجربه یاد گرفته می‌شود؟
\begin{enumerate}
  \item \term{دستور جهانی}{Universal Grammar}: انسان از بدو تولد اصول زبانیِ جهان‌شمول دارد و متناسب با محیط «پارامترها» را تنظیم می‌کند (کودک فارسی‌زبان فارسی و کودک انگلیسی‌زبان انگلیسی یاد می‌گیرد). در کلاس با تنظیم پارامترهای یک مدل مقایسه شد: ساختار اولیه هست و تجربه شکل نهایی را تعیین می‌کند.
  \item \term{یادگیری توزیعی}{Distributional / Statistical Learning}: انسان با دیدن نمونه‌های بسیار، قواعد را از الگوهای آماری داده یاد می‌گیرد.
\end{enumerate}
نقدهای واردشده به تفسیر سخت‌گیرانهٔ دستور جهانی: تفاوت‌های ساختاری میان زبان‌ها؛ شواهد یادگیری الگوهای جدید از داده؛ ابهام دربارهٔ اینکه دقیقاً کدام ویژگی‌ها فطری‌اند. \textbf{جمع‌بندی کلاس:} ایدهٔ زایشی و سلسله‌مراتبی‌بودن زبان همچنان مهم است، اما ادعای مجموعهٔ مشخصی از قواعد دستوری مادرزادی محل بحث است و میزان فطری‌بودن زبان هنوز موضوعی پژوهشی و مورد اختلاف است.

\begin{correctionbox}{دقت تاریخی دربارهٔ نظریهٔ چامسکی}
سال ۱۹۵۷ سال انتشار \eng{Syntactic Structures} است؛ نقد چامسکی بر رویکرد تقلیدی/رفتارگرا (کتاب اسکینر) در ۱۹۵۹ آمد و صورت‌بندی «اصول و پارامترها» بعدتر (دههٔ ۱۹۸۰) پدید آمد. دو دیدگاه الزاماً دو سر یک طیف‌اند: بسیاری از پژوهشگران امروز بر ترکیبی از سوگیری‌های ذاتی و یادگیری از داده تأکید می‌کنند.
\end{correctionbox}

\section{ورود به آواشناسی}

\term{آواشناسی}{Phonetics} دو بخش اصلی دارد:
\begin{description}
  \item[\term{آواشناسی تولیدی}{Articulatory Phonetics}] اینکه بدن چگونه صدا می‌سازد (لب‌ها، زبان، فک و تارهای صوتی چه می‌کنند). موضوع این جلسه.
  \item[\term{آواشناسی آکوستیکی}{Acoustic Phonetics}] خود موج صوتی و ویژگی‌های فیزیکی آن. موضوع جلسهٔ بعد.
\end{description}
هدف نهایی: پس از شناخت دادهٔ گفتاری، آن را به بازنمایی مناسب پردازش رایانه‌ای تبدیل کنیم.

واژهٔ \eng{articulation} یعنی نحوهٔ حرکت و هماهنگی اندام‌های گفتاری (\eng{articulators}: لب‌ها، زبان، فک، نرم‌کام و \ldots)؛ در کلاس برای شهود به آرتیکولاتور مکانیکی دندان‌پزشکی (شبیه‌ساز حرکت فک) اشاره شد.

\section{آناتومی تولید گفتار}

بدن را با سه صفحهٔ مرجع بررسی می‌کنند: \term{ساجیتال}{Sagittal} (نیمهٔ چپ/راست، نمای از پهلو)، \term{فرونتال}{Frontal / Coronal} (جلو/پشت) و \term{افقی}{Horizontal / Transverse} (بالا/پایین). برای مطالعهٔ حنجره، نای، دهان و مسیر هوا \textbf{نمای ساجیتال} بیشترین کاربرد را دارد (شکل \ref{fig:ch02-sagittal}).

\boardfigure{0.45}{ch02-vocal-tract-sagittal.png}{برش ساجیتال مجرای صوتی: حفرهٔ بینی، حلق، تارهای صوتی (درون حنجره)، نای و ریه‌ها.}
\label{fig:ch02-sagittal}

\begin{table}[htbp]
\centering
\caption{اندام‌های اصلی دخیل در تولید گفتار}
\small
\begin{tabular}{p{2.9cm} p{3.6cm} p{6.6cm}}
\toprule
\textbf{اندام} & \textbf{معادل انگلیسی} & \textbf{نقش} \\
\midrule
ریه‌ها & \eng{Lungs} & منبع جریان هوا (تأمین فشار) \\
دیافراگم & \eng{Diaphragm} & ماهیچهٔ اصلی دم \\
نای & \eng{Trachea} & مسیر هوا از ریه به حنجره \\
حنجره & \eng{Larynx} & محل قرارگیری تارهای صوتی \\
تارهای صوتی & \eng{Vocal Folds} & تولید صوت واک‌دار از راه ارتعاش \\
حلق & \eng{Pharynx} & بخشی از مجرای صوتی بالای حنجره \\
حفرهٔ بینی & \eng{Nasal Cavity} & تولید صداهای خیشومی \\
دهان، زبان، لب‌ها، فک، دندان‌ها & \eng{Oral Articulators} & شکل‌دهی و تمایز صداها \\
\bottomrule
\end{tabular}
\end{table}

تارهای صوتی \textbf{زیر حلق} و در ابتدای نای، درون حنجره قرار دارند. می‌توان اندام‌ها را در سه گروه دید: سامانهٔ تنفسی (انرژی و فشار)، حنجره («منبع صدا») و مجرای صوتی (فیلتر و شکل‌دهی)؛ این همان مدل «منبع–فیلتر» است که در جلسهٔ بعد می‌آید.

\subsection{دم و بازدم}
\begin{description}
  \item[دم (\eng{Inhalation})] دیافراگم منقبض می‌شود و پایین می‌رود؛ حجم قفسه زیاد و فشار ریه کم می‌شود و هوا وارد می‌شود. تارهای صوتی در تنفس معمولی باز هستند.
  \item[بازدم (\eng{Exhalation})] دیافراگم به حالت استراحت برمی‌گردد و هوا خارج می‌شود. این جریان «منبع انرژی» گفتار است: هوا هنگام عبور از حنجره تارها را به ارتعاش درمی‌آورد و صدا در مجرای بالاتر شکل می‌گیرد.
\end{description}

\begin{correctionbox}{بازدم در گفتار منفعل نیست}
«دیافراگم در بازدم ریلکس است» برای تنفس آرام درست است، اما در گفتار سامانهٔ عصبی–عضلانی دائماً مداخله می‌کند تا فشار زیر حنجره در طول جمله تقریباً ثابت بماند (حدود ۵ تا ۱۰ سانتی‌متر آب): ابتدای جمله عضلات دمی بازدم را «ترمز» و انتهای آن عضلات بازدمی هوا را بیرون می‌رانند.
\end{correctionbox}

\paragraph{پرسش کلاسی: چرا معمولاً هنگام بازدم صحبت می‌کنیم؟} در دم آرام تارها باز هستند و هوا بدون ارتعاش می‌گذرد؛ برای صدای واک‌دار تارها باید نزدیک هم باشند و جریان هوا با فشار کافی از آن‌ها بگذرد. گفتار هنگام دم غیرممکن نیست (برخی افراد، در آواز یا تمرین‌های صوتی، چنین می‌کنند)، اما الگوی غالب گفتار در همهٔ زبان‌ها استفاده از هوای بازدمی است؛ زیرا فشار را می‌توان در بازدم طولانی و دقیق کنترل کرد.

\paragraph{پرسش کلاسی: نای چرا حلقه‌حلقه است؟} نای لوله‌ای است که هوا را از حنجره به ریه می‌برد و دیوارهٔ آن \textbf{حلقه‌های غضروفی} دارد: (۱) هنگام دم فشار درون نای کمتر از بیرون می‌شود و بدون حلقه‌ها مانند نی نرم هنگام مکیدن روی هم می‌خوابید؛ (۲) فاصلهٔ میان حلقه‌ها به نای و گردن انعطاف می‌دهد. پس نای باید \textbf{همیشه باز} بماند، برخلاف تارهای صوتی که مرتب باز و بسته می‌شوند.

\section{حنجره و تارهای صوتی}

هوای بازدم به حنجره می‌رسد که جفت \term{تار صوتی}{Vocal Folds} در آن است. فضای میان دو تار \term{چاکنای}{Glottis} نام دارد؛ در تنفس آرام باز است و هنگام صدای واک‌دار تارها نزدیک می‌شوند و وارد چرخهٔ ارتعاش می‌شوند.

\begin{correctionbox}{«تار» یا «چین»؟}
«تارهای صوتی» ترجمهٔ تحت‌اللفظی \eng{vocal cords} است و تصویر نادرستی از سیم کشیده‌شده می‌دهد؛ اصطلاح علمی \eng{vocal folds} (چین‌های صوتی) است: توده‌ای از بافت ماهیچه‌ای و مخاطی که با چرخش غضروف‌های آریتنوئید به هم نزدیک یا از هم دور می‌شود. در این جزوه به‌دلیل رواج همان «تارهای صوتی» را به کار می‌بریم.
\end{correctionbox}

\subsection{چرا تارها فقط یک بار باز نمی‌شوند؟ اصل برنولی}
\paragraph{پرسش کلاسی:} وقتی هوا از میان تارها می‌گذرد، چرا همان‌طور باز نمی‌مانند و مرتب باز و بسته می‌شوند؟ پاسخ با \term{اصل برنولی}{Bernoulli's Principle} داده شد.

\begin{definitionbox}{اصل برنولی}
در جریان پایا، تراکم‌ناپذیر و بدون اصطکاک در امتداد یک خط جریان:
\[
p + \tfrac{1}{2}\rho v^2 + \rho g h = \text{ثابت}
\]
که $p$ فشار ایستا، $\rho$ چگالی، $v$ سرعت و $h$ ارتفاع است. در ارتفاع ثابت، \textbf{سرعت بیشتر یعنی فشار ایستای کمتر.}
\end{definitionbox}

(آزمایش دو بطری کلاس: اگر میان دو بطری سبک هوا بدمیم، بطری‌ها به هم نزدیک می‌شوند، چون فشار ناحیهٔ پرسرعت کمتر است.)

\paragraph{چرخهٔ ارتعاش (روایت ساده‌شدهٔ کلاس).}
\begin{enumerate}
  \item تارها نزدیک هم‌اند و مجرا بسته است؛ فشار زیر تارها بالا می‌رود.
  \item فشار تارها را از هم باز می‌کند.
  \item هوا با سرعت زیاد از شکاف باریک می‌گذرد.
  \item طبق اصل برنولی فشار در شکاف کم می‌شود و همراه با ارتجاع بافت، تارها دوباره بسته می‌شوند.
  \item چرخه تکرار می‌شود و ارتعاش دوره‌ای هوا موج صوتی می‌سازد.
\end{enumerate}
چون حرکت تکرارشونده است، \textbf{فرکانس} دارد: همان \term{فرکانس پایه}{Fundamental Frequency, F0} که به‌صورت زیر یا بم‌بودن صدا ادراک می‌شود.

\begin{correctionbox}{اصل برنولی تنها توضیح ارتعاش نیست}
توضیح کلاس مدلی ساده‌شده بود. نظریهٔ پذیرفته‌شده \term{نظریهٔ ماهیچه‌ای–ارتجاعی–آیرودینامیکی}{Myoelastic-Aerodynamic Theory} است: ماهیچه‌ها تارها را نزدیک می‌کنند، فشار زیرحنجره‌ای آن‌ها را باز می‌کند و نیروی ارتجاعی بافت به‌علاوهٔ افت فشار برنولی آن‌ها را می‌بندد. برای تداوم ارتعاش، تفاوت فاز میان لبهٔ پایینی و بالایی تارها انرژی را از جریان هوا به بافت منتقل می‌کند.
\end{correctionbox}

\subsection{فرکانس ارتعاش و موج حنجره‌ای}
در هر چرخه تارها یک بار باز و بسته می‌شوند و هوا به‌صورت پالس‌های متناوب وارد مجرای صوتی می‌شود (شکل \ref{fig:ch02-glottal}). تعداد چرخه‌ها در ثانیه همان $\chtwoFz $ و دورهٔ تناوب $\chtwoTz  = 1/\chtwoFz $ است.

\coursefigure{0.55}{ch02-glottal-airflow.png}{جریان هوای عبوری از چاکنای برای $\chtwoFz  = 125$ هرتز؛ تارها هر ۸ میلی‌ثانیه یک بار باز و بسته می‌شوند. (شکل افزوده)}
\label{fig:ch02-glottal}

مثال: $\chtwoFz  = 100$ هرتز (صدای مردانه) $\Rightarrow$ $\chtwoTz  = 10$ میلی‌ثانیه؛ $\chtwoFz  = 200$ هرتز (زنانه) $\Rightarrow$ $5$ میلی‌ثانیه. $\chtwoFz $ مردان بزرگسال تقریباً ۸۵ تا ۱۸۰ و زنان ۱۶۵ تا ۲۵۵ هرتز است؛ تارهای کوتاه‌تر و نازک‌تر فرکانس بالاتری دارند.

\section{زنجیرهٔ گفتار}

یک مفهوم ذهنی چگونه به فرمان عضلانی، صدا و جمله‌ای قابل‌فهم برای مخاطب تبدیل می‌شود؟ مدل \term{زنجیرهٔ گفتار}{Speech Chain} (دنس و پینسون، ۱۹۶۳) این مسیر را از مغز گوینده تا مغز شنونده نشان می‌دهد (شکل \ref{fig:ch02-chain-tikz}).

\begin{figure}[htbp]
\centering
\begin{tikzpicture}[rtl,
  rtl,
  node distance=4mm,
  box/.style={draw, rounded corners, align=center, minimum height=1.15cm, text width=1.95cm, font=\scriptsize, fill=softblue},
  wbox/.style={box, fill=softyellow},
  fbox/.style={box, fill=softgreen, text width=3.4cm},
  arr/.style={-{Stealth[length=2mm]}, thick},
  lab/.style={font=\scriptsize\bfseries, text=mainblue}]
\node[box] (b1) {مغز گوینده\\ (قصد و کدگذاری زبانی)};
\node[box, left=of b1] (m) {نورون‌های حرکتی\\ و ماهیچه‌ها};
\node[box, left=of m] (a) {اندام‌های گفتار\\ \lr{(Articulators)}};
\node[wbox, left=of a] (w) {موج صوتی\\ \lr{(Sound waves)}};
\node[box, left=of w] (e) {گوش و نورون‌های حسی شنونده};
\node[box, left=of e] (b2) {مغز شنونده\\ (فهم زبانی)};
\draw[arr] (b1) -- (m); \draw[arr] (m) -- (a); \draw[arr] (a) -- (w); \draw[arr] (w) -- (e); \draw[arr] (e) -- (b2);
\node[fbox, below=9mm of m] (f) {گوش و نورون‌های حسی گوینده\\ (بازخورد شنیداری)};
\draw[arr, dashed] (w.south) |- (f.west);
\draw[arr, dashed] (f.east) -| (b1.south);
\node[lab, above=1mm of b1] {سطح زبانی};
\node[lab, above=1mm of m] {سطح فیزیولوژیک};
\node[lab, above=1mm of w] {سطح آکوستیک};
\node[lab, above=1mm of e] {سطح فیزیولوژیک};
\node[lab, above=1mm of b2] {سطح زبانی};
\end{tikzpicture}
\caption{زنجیرهٔ گفتار با سه سطح زبانی، فیزیولوژیک و آکوستیک. خط‌چین: بازخورد شنیداری گوینده. (شکل افزوده)}
\label{fig:ch02-chain-tikz}
\end{figure}

مراحل: مغز گوینده (قصد بیان) $\leftarrow$ \term{نورون‌های حرکتی}{Motor Neurons} (فرمان به ماهیچه‌ها) $\leftarrow$ حرکت هماهنگ اندام‌های گفتار $\leftarrow$ موج صوتی در هوا $\leftarrow$ گوش و \term{نورون‌های حسی}{Sensory Neurons} (تبدیل ارتعاش هوا به سیگنال عصبی) $\leftarrow$ مغز شنونده. افزون بر این، \textbf{گوینده صدای خود را هم می‌شنود}؛ پس یک حلقهٔ \term{بازخورد}{Feedback} شنیداری برای کنترل و اصلاح گفتار وجود دارد.

\subsection{هم‌تولیدی}
در کلاس ادای کلمهٔ «سیب» به نواختن پیانو تشبیه شد (هر کلید یک واج)؛ اما این تشبیه محدودیت دارد: واج‌ها یکی‌یکی و مجزا تولید نمی‌شوند. حرکات واج‌های مجاور تا حدی \textbf{هم‌پوشانی} دارند و اندام‌ها در حین یک واج برای واج بعدی آماده می‌شوند. این پدیده \term{هم‌تولیدی}{Coarticulation} نام دارد. مثلاً در «سو» (\lr{/su/}) لب‌ها از همان \lr{/s/} گرد می‌شوند، ولی در «سی» (\lr{/si/}) کشیده‌اند.

\begin{exambox}{اهمیت هم‌تولیدی در پردازش گفتار}
گفتار فرایندی پیوسته و هماهنگ است، نه توالی حرکات مجزا؛ پس نمی‌توان فرض کرد که در سیگنال صوتی همهٔ واج‌ها مرزهای کاملاً مشخص و مستقل دارند.
\end{exambox}

\section{واج، همخوان، واکه و هجا}

\begin{definitionbox}{واج}
\term{واج}{Phoneme} یکی از کوچک‌ترین واحدهای تمایزدهندهٔ معنایی در نظام آوایی یک زبان است: با جایگزینی یک واج در واژه ممکن است معنا عوض شود. مثال: «بار» و «پار» (\lr{/b/} و \lr{/p/})، «سیر» و «زیر» (\lr{/s/} و \lr{/z/}) \term{جفت کمینه}{Minimal Pair} هستند.
\end{definitionbox}

واج‌ها دو دستهٔ اصلی‌اند: \term{همخوان (صامت)}{Consonant} و \term{واکه (مصوت)}{Vowel}. از کنار هم قرار گرفتن واج‌ها، \term{هجا}{Syllable} و سپس کلمه ساخته می‌شود. الگوی هجا را با \eng{C} (همخوان) و \eng{V} (واکه) نشان می‌دهند؛ مثلاً «کتاب» دو هجا دارد: «کِ» (\eng{CV}) و «تاب» (\eng{CVC}). خود صدای فیزیکی را \term{آوا}{Phone} می‌گویند و واج رده‌ای انتزاعی از آواهای معادل است.

\section{واکه‌ها و نقشهٔ جایگاه زبان}

زبان (و فک و لب‌ها) برای هر واکه در موقعیت مشخصی از دهان قرار می‌گیرد؛ این موقعیت‌ها را در \term{نقشهٔ واکه‌ها}{Vowel Chart} نمایش می‌دهند. دو بُعد اصلی:
\begin{enumerate}
  \item \textbf{جلو یا عقب بودن زبان} (\eng{Front / Central / Back});
  \item \textbf{ارتفاع زبان} (بالا–میانه–پایین؛ در نمودار بسته–میانه–باز).
\end{enumerate}
تعداد واکه‌ها در زبان‌ها فرق می‌کند: نمودار انگلیسی کلاس ۱۲ واکه داشت، در حالی که در \textbf{فارسی معیار شش واکهٔ اصلی} در نظر می‌گیریم.

\begin{table}[htbp]
\centering
\caption{واکه‌های فارسی معیار و جایگاه تقریبی زبان}
\begin{tabular}{c c c c l}
\toprule
\textbf{واکه} & \textbf{نماد \lr{IPA}} & \textbf{نمونه} & \textbf{جلو/عقب} & \textbf{ارتفاع زبان} \\
\midrule
ای & \lr{/i/} & دیر & جلو & بالا (بسته) \\
اِ & \lr{/e/} & دل & جلو & میانه \\
اَ & \lr{/æ/} & در & جلو & پایین (باز) \\
آ & \lr{/ɒ/} یا \lr{/ɑ/} & آب & عقب & پایین (باز) \\
اُ & \lr{/o/} & گل & عقب & میانه \\
او & \lr{/u/} & دور & عقب & بالا (بسته) \\
\bottomrule
\end{tabular}
\end{table}

\begin{figure}[htbp]
\centering
\begin{tikzpicture}[rtl, scale=0.95,
  v/.style={circle, fill=mainblue, inner sep=1.6pt},
  lbl/.style={font=\small, fill=white, inner sep=1pt}]
\draw[thick] (0,3) -- (4,3) -- (4,0) -- (1.6,0) -- cycle;
\draw[dashed, gray] (0.55,2) -- (4,2);
\draw[dashed, gray] (1.1,1) -- (4,1);
\node[v] (i) at (0.1,3) {};  \node[lbl, below right=1mm and 1mm of i] {\lr{/i/} ای};
\node[v] (e) at (0.55,2) {}; \node[lbl, below right=1mm and 1mm of e] {\lr{/e/} اِ};
\node[v] (ae) at (1.2,0.75) {}; \node[lbl, right=2mm of ae] {\lr{/æ/} اَ};
\node[v] (u) at (3.9,3) {}; \node[lbl, below left=1mm and 1mm of u] {\lr{/u/} او};
\node[v] (o) at (3.9,2) {}; \node[lbl, below left=1mm and 1mm of o] {\lr{/o/} اُ};
\node[v] (a) at (3.9,0) {}; \node[lbl, above left=1mm and 1mm of a] {\lr{/ɒ/} آ};
\node[font=\footnotesize] at (0.2,3.6) {جلو};
\node[font=\footnotesize] at (3.8,3.6) {عقب};
\node[font=\footnotesize, anchor=east] at (-0.2,3) {بسته};
\node[font=\footnotesize, anchor=east] at (0.3,1.5) {میانه};
\node[font=\footnotesize, anchor=east] at (1.4,-0.3) {باز};
\end{tikzpicture}
\caption{نقشهٔ تقریبی واکه‌های فارسی روی ذوزنقهٔ واکه‌ها (جلو در سمت چپ، مطابق قرارداد \lr{IPA}). (شکل افزوده)}
\label{fig:ch02-vowel-persian}
\end{figure}

\begin{assignmentbox}{تمرین کلاسی: رسم نقشهٔ واکه‌های فارسی}
جایگاه شش واکهٔ فارسی را روی نمودار واکه‌ها مشخص کنید. هر واکه را تلفظ کنید و به حرکت زبان توجه کنید: «او» (زبان بالا و عقب) و سپس «آ» (زبان پایین و عقب‌تر)، و همین‌طور «ای» و «اَ». دربارهٔ جایگاه «آ»، «اَ» و «او» در کلاس بحث شد. هدف: درک تجربی این نکته که واکه‌ها با تغییر شکل و جایگاه اندام‌های گفتار از هم متمایز می‌شوند. (شکل \ref{fig:ch02-vowel-persian} پاسخ تقریبی است.)
\end{assignmentbox}

بُعد سوم توصیف واکه‌ها \textbf{گردی لب‌ها} است (در فارسی \lr{/u/}، \lr{/o/} و \lr{/ɒ/} گرد هستند). در جلسهٔ بعد می‌بینیم که ارتفاع زبان با فرمنت اول ($F_1$) و جلو/عقب‌بودن با فرمنت دوم ($F_2$) ارتباط مستقیم دارد.

\section{همخوان‌ها و ویژگی‌های تولیدی آن‌ها}

همخوان‌ها را با سه ویژگی توصیف می‌کنیم: \term{محل تولید}{Place of Articulation}، \term{شیوهٔ تولید}{Manner of Articulation} و \term{واک‌داری}{Voicing}. در کلاس هر ویژگی با یک تمرین عملی همراه بود.

\subsection{محل تولید}
محل تولید مشخص می‌کند مسدودشدن یا تنگ‌شدن جریان هوا در کدام بخش دستگاه گفتار رخ می‌دهد:
\begin{itemize}
  \item «ب» و «پ»: دو لب $\leftarrow$ \term{دولبی}{Bilabial};
  \item «ف» و «و» (\lr{/v/}): لب پایین و دندان‌های بالا $\leftarrow$ \term{لبی‌دندانی}{Labiodental};
  \item «س» و «ز»: زبان نزدیک لثه $\leftarrow$ \term{لثوی}{Alveolar};
  \item «ه»: ناحیهٔ چاکنای $\leftarrow$ \term{چاکنایی}{Glottal}.
\end{itemize}
(در فارسی معیار «و» همخوان \lr{/v/} است نه \lr{/w/}، و «ح» و «ه» هر دو \lr{/h/} تلفظ می‌شوند.)

\subsection{شیوهٔ تولید}
شیوهٔ تولید مشخص می‌کند جریان هوا چگونه کنترل می‌شود:
\begin{description}
  \item[\term{انسدادی}{Stop / Plosive}] جریان هوا لحظه‌ای کاملاً بسته و سپس رها می‌شود (ضربه‌ای کوچک روی دست جلوی دهان هنگام «پ…پ…پ»). مثال: پ، ب، ت، د، ک، گ.
  \item[\term{سایشی}{Fricative}] هوا ناگهانی رها نمی‌شود؛ از مجرایی باریک می‌گذرد و صدای سایشی می‌سازد («ففف…»). مثال: ف، س، ز، ش.
  \item[\term{خیشومی}{Nasal}] هوا از بینی می‌گذرد («مممم…»). مثال: م، ن.
\end{description}
در خیشومی‌ها \term{نرم‌کام}{Velum} پایین می‌آید و مسیر بینی باز می‌شود، اما در دهان انسداد کامل هم هست (برای \lr{/m/} دو لب)؛ اگر بینی را بگیرید «م» تقریباً تولید نمی‌شود ولی «ب» می‌شود.

فارسی معیار حدود ۲۳ همخوان دارد. جدول زیر آن‌ها را بر اساس محل و شیوه دسته‌بندی می‌کند؛ در هر خانه ابتدا بی‌واک و سپس واک‌دار می‌آید.

\begin{table}[htbp]
\centering
\caption{همخوان‌های فارسی بر حسب محل، شیوه و واک‌داری}
\footnotesize
\setlength{\tabcolsep}{3.5pt}
\begin{tabular}{l c c c c c c c c}
\toprule
\textbf{شیوه} \textbackslash\ \textbf{محل} & دولبی & لبی‌دندانی & لثوی & پس‌لثوی & کامی & نرم‌کامی & ملازی & چاکنایی \\
\midrule
انسدادی & \lr{p\ b} & & \lr{t\ d} & & & \lr{k\ ɡ} & \lr{ɢ} & \lr{ʔ} \\
انسدادی–سایشی & & & & \lr{tʃ\ dʒ} & & & & \\
سایشی & & \lr{f\ v} & \lr{s\ z} & \lr{ʃ\ ʒ} & & \lr{x} & & \lr{h} \\
خیشومی & \lr{m} & & \lr{n} & & & & & \\
روان کناری & & & \lr{l} & & & & & \\
زنشی & & & \lr{ɾ} & & & & & \\
غلت (نیم‌واکه) & & & & & \lr{j} & & & \\
\bottomrule
\end{tabular}
\end{table}

\subsection{واک‌داری}
\begin{definitionbox}{واک‌داری}
در صداهای \term{واک‌دار}{Voiced} تارهای صوتی ارتعاش دوره‌ای دارند (ز، ب، د، م و همهٔ واکه‌ها)؛ در صداهای \term{بی‌واک}{Voiceless} چنین ارتعاشی نیست (س، پ، ت، ف).
\end{definitionbox}

\begin{assignmentbox}{تمرین کلاسی: «س» و «ز»}
دست خود را روی گلو بگذارید و «سسسس» و سپس «زززز» را کشیده بگویید. هنگام «ز» ارتعاش تارها را حس می‌کنید، هنگام «س» نه. همین را برای «ب» و «پ» هم انجام دهید.
\end{assignmentbox}

\begin{table}[htbp]
\centering
\caption{جفت‌هایی که فقط در واک‌داری تفاوت دارند}
\begin{tabular}{l c c c c}
\toprule
\textbf{ویژگی} & \textbf{\lr{/s/} س} & \textbf{\lr{/z/} ز} & \textbf{\lr{/p/} پ} & \textbf{\lr{/b/} ب} \\
\midrule
محل تولید & لثوی & لثوی & دولبی & دولبی \\
شیوهٔ تولید & سایشی & سایشی & انسدادی & انسدادی \\
واک‌داری & بی‌واک & واک‌دار & بی‌واک & واک‌دار \\
\bottomrule
\end{tabular}
\end{table}

پس دو همخوان می‌توانند در محل و شیوه یکسان باشند و فقط در واک‌داری فرق کنند؛ این برای خطاهای بازشناسی اهمیت ویژه دارد (بخش بعد).

\begin{correctionbox}{واک‌داری و «ارتعاش» همیشه یکی نیستند}
تعریف «واک‌دار = ارتعاش تارها» ساده‌شده است: تمایز \lr{/p/} و \lr{/b/} بیشتر با \term{زمان آغاز واک}{Voice Onset Time, VOT} و دمیدگی مشخص می‌شود و گاهی \lr{/b/} در آغاز واژه کاملاً واک‌دار نیست. برای شروع، همان تعریف ارتعاش کافی است.
\end{correctionbox}

\section{کاربرد آواشناسی در هوش مصنوعی}

پرسش استاد: «این همه مطلب دربارهٔ زبان، تارهای صوتی و آواشناسی به چه درد ما به‌عنوان مهندس هوش مصنوعی می‌خورد؟» کاربردهای مطرح‌شده:
\begin{itemize}
  \item \textbf{مدل‌های ساده‌تر:} در \eng{STT} مبتنی بر یادگیری عمیق لازم نیست ویژگی‌های آوایی را دستی وارد کنیم، اما برای هر مسئله به مدل بزرگ نیاز نیست؛ اگر ثابت کنیم مثلاً تنها با \eng{Place of Articulation} کلاس‌ها تفکیک‌پذیرند، یک طبقه‌بند ساده کافی است (شرط: ابتدا باید توان تفکیک ویژگی‌ها بررسی و اثبات شود).
  \item \textbf{کمبود داده (\eng{Low-Resource}):} محل، شیوه و واک‌داری دانش مفیدی دربارهٔ ساختار داده می‌دهند تا مدل با داده کمتر به عملکرد مناسب برسد.
  \item \textbf{افزودن دانش (\eng{Knowledge Injection}):} ویژگی‌های آواشناسی را به‌صورت برچسب یا ویژگی به داده می‌افزاییم؛ مثلاً برای \lr{/p/}: \eng{Bilabial}، \eng{Stop}، \eng{Voiceless}.
  \item \textbf{مدل‌های سبک‌تر:} اگر ویژگی‌ای را از قبل می‌شناسیم، می‌توان آن را مستقیم به مراحل بعدی داد (شبیه کاهش ابعاد مسئله). صرف افزودن ویژگی مدل را خودکار کوچک نمی‌کند؛ معماری و آموزش هم باید متناسب طراحی شوند.
  \item \textbf{تحلیل خطا و تفسیرپذیری} (مهم‌ترین کاربرد به نظر استاد).
\end{itemize}

\begin{exambox}{تفاوت \eng{Knowledge Injection} و \eng{Data Cleaning}}
در \term{پاک‌سازی داده}{Data Cleaning} داده‌های نامناسب یا خراب را اصلاح می‌کنیم؛ اما در \term{تزریق دانش}{Knowledge Injection} اطلاعاتی را که دانش پیشین ما دربارهٔ مسئله است به داده یا فرایند آموزش می‌افزاییم. این دو یکی نیستند.
\end{exambox}

\subsection{تحلیل خطا و تفسیرپذیری}
فرض کنید برای یک مدل \eng{STT} نرخ خطای کلمه (\eng{WER}) برابر ۱۲ درصد گزارش شده است. این عدد فقط می‌گوید مدل خطا دارد، نه اینکه چه نوع خطاهایی بیشتر رخ می‌دهد. برای تحلیل دقیق‌تر، خروجی مدل (\eng{Prediction}) را با مرجع صحیح (\eng{Ground Truth}) مقایسه و ویژگی‌های واجیِ واج‌های اشتباه‌شده را بررسی می‌کنیم: از نظر محل، شیوه یا واک‌داری چه شباهت یا تفاوتی دارند؟ مثلاً اگر مدل «ب» و «پ» را زیاد با هم اشتباه بگیرد:

\begin{center}
\begin{tabular}{l c c}
\toprule
\textbf{ویژگی} & \lr{/b/} & \lr{/p/} \\
\midrule
\eng{Place} & \eng{Bilabial} & \eng{Bilabial} \\
\eng{Manner} & \eng{Stop} & \eng{Stop} \\
\eng{Voicing} & \eng{Voiced} & \eng{Voiceless} \\
\bottomrule
\end{tabular}
\end{center}
محل و شیوه یکی است و تفاوت اصلی در واک‌داری است. به‌جای «۱۲ درصد خطا» می‌توان گفت چه سهمی از خطاها به واک‌داری و چه سهمی به محل تولید مربوط است. \textbf{بخش قابل‌توجهی از خطاها بین واج‌هایی رخ می‌دهد که از نظر ویژگی‌های آرتیکولاتوری مشابه‌اند.} پس از شناسایی الگوی خطا می‌توان داده‌های برخی واج‌ها را زیاد کرد، کیفیت و تنوع داده را بهبود داد، ویژگی‌های مرتبط را افزود یا معماری و آموزش را تغییر داد.

\begin{addedbox}{مثال عددی: \eng{WER} و فاصلهٔ ویژگی‌ها}
$\text{\lr{WER}} = (S+D+I)/N$ (جایگزینی، حذف، درج؛ تقسیم بر تعداد کلمات مرجع): برای $N=50$، $S=4$، $D=1$، $I=1$ داریم $6/50 = 12\%$. اگر هر همخوان را با بردار دودویی (دولبی، لثوی، انسدادی، سایشی، واک‌دار) نمایش دهیم، $\text{\lr{b}}=(1,0,1,0,1)$، $\text{\lr{p}}=(1,0,1,0,0)$ و $\text{\lr{s}}=(0,1,0,1,0)$؛ فاصلهٔ همینگ $d(b,p)=1$ اما $d(b,s)=5$: اشتباه «ب» با «پ» خطایی نزدیک و قابل‌انتظار و اشتباه «ب» با «س» جدی‌تر است.
\end{addedbox}

\section{جمع‌بندی}
در این فصل رویکرد پایین‌به‌بالا و جایگاه پردازش گفتار در علوم زبان را دیدیم، نظریهٔ زایشی چامسکی را مرور کردیم و سپس به آواشناسی تولیدی رسیدیم: آناتومی دستگاه گفتار، تنفس، ارتعاش تارهای صوتی، زنجیرهٔ گفتار و هم‌تولیدی، و توصیف واکه‌ها و همخوان‌ها. در پایان دیدیم این دانش در مدل‌های ساده‌تر، تزریق دانش، کمبود داده و به‌ویژه تحلیل خطا به کار می‌آید. فصل بعد به آواشناسی آکوستیکی می‌پردازد.

\begin{exambox}{نکات مهم برای مرور و امتحان}
\begin{itemize}
  \item رویکرد جلسه \textbf{پایین‌به‌بالا} است: ابتدا ماهیت داده، سپس بازنمایی، سپس مدل. گفتار سیگنالی \textbf{فیزیکی و پیوسته} است؛ متن \textbf{نمادین و گسسته}.
  \item درس ما در چارچوب \textbf{زبان‌شناسی رایانشی} است؛ هدف، شناخت مفاهیم زبان‌شناسی برای استفاده در مدل‌هاست.
  \item چامسکی: زبان زایشی و سلسله‌مراتبی است؛ بحث \textbf{دستور جهانی} در برابر \textbf{یادگیری آماری/توزیعی} هنوز باز است.
  \item مسیر تولید صدا: ریه (بازدم) $\leftarrow$ نای $\leftarrow$ حنجره و تارهای صوتی $\leftarrow$ مجرای صوتی (حلق، دهان، بینی).
  \item تارهای صوتی با فشار هوا، ارتجاع بافت و اثر \textbf{برنولی} (سرعت بیشتر، فشار کمتر) دوره‌ای باز و بسته می‌شوند؛ $\chtwoTz  = 1/\chtwoFz $. نای برای باز ماندن حلقه‌دار است، ولی تارها باید باز و بسته شوند.
  \item \textbf{زنجیرهٔ گفتار:} مغز گوینده $\leftarrow$ نورون حرکتی $\leftarrow$ اندام‌ها $\leftarrow$ موج صوتی $\leftarrow$ گوش و نورون حسی $\leftarrow$ مغز شنونده، به‌علاوهٔ بازخورد شنیداری.
  \item \textbf{هم‌تولیدی:} هم‌پوشانی حرکات واج‌های مجاور؛ مرز واج‌ها در سیگنال کاملاً مشخص نیست.
  \item واکه‌ها: جایگاه زبان (جلو/عقب و بالا/پایین)؛ فارسی شش واکه دارد (ای، اِ، اَ، آ، اُ، او).
  \item همخوان‌ها: \textbf{\eng{Place}} (دولبی، لبی‌دندانی، لثوی، چاکنایی، \ldots)، \textbf{\eng{Manner}} (انسدادی، سایشی، خیشومی، \ldots) و \textbf{\eng{Voicing}}. مثال‌های کلیدی: «ب/پ» و «س/ز».
  \item کاربرد در هوش مصنوعی: مدل‌های ساده‌تر، کمبود داده، \eng{Knowledge Injection} (که با \eng{Data Cleaning} فرق دارد)، مدل‌های سبک‌تر و به‌ویژه تحلیل خطا: به‌جای یک عدد کلی مثل \eng{WER}، بررسی الگو و علت خطاها با ویژگی‌های واجی.
\end{itemize}
\end{exambox}
